% ============================================================
%  PARTE II — ARQUETA DE FANGOS MEZCLA
% ============================================================
\chapter{Arqueta de fangos mezcla}
\label{chap:arqueta}

Se proyecta una nueva arqueta de fangos mezcla con las siguientes características:

\begin{table}[H]
  \centering
  \caption{Parámetros de diseño de la arqueta de fangos mezcla}
  \label{tab:arqueta_datos}
  \begin{tabular}{@{}lcc@{}}
    \toprule
    \textbf{Parámetro} & \textbf{Símbolo} & \textbf{Valor} \\
    \midrule
    Largo interior                         & $L$                 & \largoArqueta~m \\
    Ancho interior                         & $B$                 & \anchoArqueta~m \\
    Altura interior                        & $H$                 & \alturaArqueta~m \\
    Espesor uniforme (paredes, losa, tapa) & $e$                 & 0,30~m \\
    Peso específico del fango              & $\gamma_f$          & 1,10~Tn/m³ \\
    Peso específico del hormigón armado    & $\gamma_{HA}$       & 2,50~Tn/m³ \\
    Sulfatos suelo                         & SO$_4^{2-}$         & 3\,500~mg/kg \\
    Clase de exposición                    & EHE-08              & IIb+Qc \\
    \bottomrule
  \end{tabular}
\end{table}

Las dimensiones exteriores resultan $L_\text{ext}=4{,}60$~m,
$B_\text{ext}=3{,}60$~m, con una altura exterior total de $4{,}10$~m
(losa + interior + tapa).
La cimentación se selecciona atendiendo a criterios técnicos: capacidad
portante, asientos, posición del nivel freático y condicionantes
medioambientales de la zona protegida.
